\section*{The Propers: Themes and Movement}

\begin{massbox}{The Lord assumes the debt before the steward answers}
Mercy does not evade justice by manipulating the account. St.~Peter Chrysologus reads Christ as the rich Lord and Judge who assumes and pays the human steward's former debt; released by that costly mercy, Spirit-led heirs are summoned to repentance, truthful administration, humble offering, and embodied restoration.
\end{massbox}

\begin{center}
\begin{tikzpicture}[x=3.55cm]
\node[guide node,text width=2.65cm] (mercy) at (0,0) {\textbf{1. Mercy precedes}\\Gift --- account};
\node[guide node,text width=2.65cm] (heirs) at (1,0) {\textbf{2. Debtors become heirs}\\Spirit --- freedom};
\node[guide emphasis,text width=2.65cm] (repent) at (2,0) {\textbf{3. Christ pays}\\Repent --- offer};
\node[guide node,text width=2.65cm] (restore) at (3,0) {\textbf{4. Communion restores}\\Mind --- body};
\foreach \x/\y in {mercy/heirs,heirs/repent,repent/restore}
  \draw[guide arrow] (\x) -- (\y);
\end{tikzpicture}
\end{center}

\noindent The four stages are a reading of the appointed liturgical sequence, and each rests on the sources cited within it. The cited Fathers expound particular texts, not the complete Roman formulary.

\subsection*{Mercy precedes the account\enspace\properrefs{Int., Coll., Gosp.}}

The entrance says \latin{Suscepimus}: mercy has been received before any human account is rendered. Yet the same antiphon confesses God's right hand full of justice, and the Collect asks him for the spirit to think and do what is \latin{recta}. Chrysologus hears the same order within the Gospel. The master requests the account now so that repentance may occur before final judgment; Christ, the rich Lord, has already assumed the steward and his former debts.

\begin{studybox}{Chrysologus: Christ the rich Lord, not the steward}
In Chrysologus, Christ is the rich man, master, and Judge; the steward is first humanity and then, in Sermon 126, the converted Gentile people. Christ says that he assumed and paid the prior debts, entered judgment though Judge, endured penalties though innocent, accepted death, and entered the underworld to raise humanity to his dignity (\emph{Sermons} 125.3--8). He is never identified with the steward's injustice.
\end{studybox}

\sectionguard
\subsection*{Debtors to the flesh become Spirit-led heirs\enspace\properrefs{Ep., Grad., All., Gosp.}}

Paul begins \latin{debitores sumus, non carni}; Luke speaks of the master's \latin{debitores}. The obligations are not identical: Romans denies the flesh's mastery, while the parable narrates liabilities within an estate. Their liturgical juxtaposition nevertheless clarifies what mercy does. It does not abolish creaturely answerability; it frees persons from slavery so that the Spirit may lead them as adopted children and coheirs.

\begin{threestable}{Relation}{Witness}{Distinction}
\textbf{Death and life} & Chrysostom treats \latin{moriemini} as a real present and final warning, while mortification concerns evil deeds, not bodily substance. & The body remains God's good work and is ordered to restoration.\\
\textbf{Adoption and refuge} & Augustine and Thomas distinguish servile terror from filial fear; the Gradual directs confidence away from self to God. & Adoption preserves reverence and judgment; Romans 9:4 also names Israel's adoption.\\
\textbf{Inheritance and suffering} & The missal stops at \latin{coheredes autem Christi}; canonical verse 17 continues with suffering with Christ unto shared glory. & Suffering does not earn adoption or become good in itself; filial conformity remains cruciform.\\
\end{threestable}
